% =====================================================================
% 4. Data
% =====================================================================
\section{Data}
\label{sec:data}

\subsection{MARKET DATA}
The market file contains daily levels of the SPI and of the SPX. Five features of the raw data required a decision.
\begin{itemize}
\item \emph{Type of index.} The SPX series is at about $13{,}000$ points in January 2025, which is consistent with the S\&P~500 \emph{total return} index and not with the price index (below $6{,}000$ at that date). Since the SPI is a total-return index by construction, both series include dividends, which is the consistent choice for asset-value returns.
\item \emph{Start of the sample.} The file starts in 1987/88 although the assignment states 1 January 2000. We use the sample from 1 January 2000, as stated in the assignment. The SPI sheet opens with twelve monthly values for 1987 before the daily series begins on 4 January 1988; they are not daily observations and are discarded in any case.
\item \emph{End of the sample.} The series end on \val{q4.raw_last_SPX|date} (SPX) and \val{q4.raw_last_SPI|date} (SPI), not on 31 October 2026 as stated in the assignment. We use all common observations, the last one being \val{q4.last_day|date}.
\item \emph{Calendars.} The two markets have different holiday calendars. We keep only the \val{q4.n_days|int} days on which both indices are quoted, so that no return combines a fresh price with a stale one.
\item \emph{Asynchronous closing.} The Swiss market closes several hours before the US market. This biases daily correlations downwards, but is of minor importance at the weekly frequency.
\end{itemize}
Weekly log-returns are computed from Wednesday close to Wednesday close (the last common trading day before it when Wednesday is a holiday), and are \emph{non-overlapping}, which gives $n=\val{q4.n_weeks|int}$ observations from \val{q4.first_week|date} to \val{q4.last_week|date}. Wednesday is the weekday least affected by holidays and by weekend effects. Overlapping returns would multiply the number of observations by five but induce strong serial dependence, which contradicts the i.i.d.\ assumption behind the likelihoods of Section~\ref{sec:factormodels}.

\subsection{STYLISED FACTS}
\begin{table}[tb]
\caption{Weekly log-returns of SPI and SPX, \val{q4.first_week|date} to \val{q4.last_week|date} ($n=\val{q4.n_weeks|int}$).}
\label{tab:desc}
\centering
\qrmtable{tab_descriptive}
\end{table}

Table~\ref{tab:desc} shows that both series are left-skewed and strongly leptokurtic; the Jarque--Bera statistics exceed any conventional critical value. (The largest weekly fall and the largest weekly rise of the SPI, in the weeks ending \val{q4.spi_min_week|date} and \val{q4.spi_max_week|date}, happen to have almost the same size.) The linear correlation between the two indices is \val{q4.corr_all|f3}. Extreme losses also occur jointly more often than a Gaussian copula allows: the probability that SPX is below its empirical $1\%$ quantile given that SPI is below its own is \val{q4.joint_exceed_emp.0.01|pct0}\% in the data, against \val{q4.joint_exceed_gauss.0.01|pct0}\% for a Gaussian copula with the same correlation (at the $5\%$ level: \val{q4.joint_exceed_emp.0.05|pct0}\% against \val{q4.joint_exceed_gauss.0.05|pct0}\%). Figure~\ref{fig:qq} confirms that the departure from normality is concentrated in the left tail. The data thus depart from the Gaussian model in two ways, in the margins and in the dependence between them, which is what the comparison of M1, M2 and M3 is designed to separate.

\begin{figure}[tb]
\centering
\qrmfigure{fig_qq}
\caption{Normal QQ-plots of the standardised weekly log-returns.}
\label{fig:qq}
\end{figure}

\subsection{PORTFOLIO}
The portfolio file lists, for each of the $100$ counterparties, the rating, the exposure $E_k$, the recovery rate $R_k$, the default probability $\pi_k$, the systematic weight $\lambda_k$, the loadings $a_k$, and the canton and coordinates of its location. Total exposure is USD~\val{q5.total_exposure|usd}; \val{q5.share_exposure_inv_grade|pct0}\% of it is rated BBB or better. Default probabilities range from $0.01\%$ (AAA) to $20\%$ (CCC), recovery rates from \val{q5.R_min|f2} to \val{q5.R_max|f2} with a mean of \val{q5.R_mean|f2}, and systematic weights from \val{q5.lam_min|f2} to \val{q5.lam_max|f2} with a mean of \val{q5.lam_mean|f2}. All loadings are non-negative. The maximum possible loss $\sum_kE_k(1-R_k)$ is USD~\val{q5.sum_w|usd} and the expected loss is USD~\val{q5.EL|usd}, of which \val{q5.EL_share_CCC|pct0}\% comes from the eight CCC names and \val{q5.EL_share_B_CCC|pct0}\% from the B and CCC names together. The counterparties are located in \val{q5.n_cantons|int} Swiss cantons, with \val{q5.share_exposure_top3_cantons|pct0}\% of the exposure in the three largest (Zurich, Vaud and Bern); this geography matters only for the climate overlay of Section~\ref{sec:climate}.
